\section{Introduction}

Lattice calculations of parton distribution functions (PDFs) are now a subject of considerable interest and efforts
(see Ref.~\cite{Cichy:2018mum} for a recent review and references to extensive literature).
Modern efforts aim at the extractions of PDFs $f(x)$ themselves rather than
their $x^N$ moments.
On the lattice, this
may be achieved by switching from
local operators to current-current correlators \cite{Detmold:2005gg}.
A further idea is to start
with equal-time correlators
\mbox{\cite{Braun:2007wv}}.

X. Ji, in the paper \cite{Ji:2013dva} that strongly stimulated further development,
made a ground-breaking proposal to consider equal-time versions of nonlocal operators
defining PDFs, distribution amplitudes, generalized parton distributions, and
transverse momentum dependent distributions.
In the case of usual PDFs, the
basic concept
of Ji's approach
is a ``parton quasi-distribution'' (quasi-PDF) $Q(y,p_3)$ \cite{Ji:2013dva,Ji:2014gla},
and
PDFs are obtained from the large-momentum $p_3 \to \infty$ limit
of quasi-PDFs.

Other approaches, such as
the ``good lattice cross sections''
\cite{Ma:2014jla,Ma:2017pxb}, the Ioffe-time analysis of equal-time correlators \cite{Braun:2007wv,Bali:2017gfr,Bali:2018spj}
and the pseudo-PDF approach \cite{Radyushkin:2017cyf,Radyushkin:2017sfi,Orginos:2017kos}
are coordinate-space oriented, and extract parton distributions taking the
short-distance $z_3 \to 0$ limit.

Both the $p_3\to \infty$ and $z_3\to 0$ limits are singular,
and one needs to use {\it matching relations}
to
convert the Euclidean lattice data into
the usual light-cone PDFs. In the quasi-PDF approach, such relations were studied
for quark \mbox{\cite{Ji:2013dva,Xiong:2013bka,Ji:2015jwa,Izubuchi:2018srq}} and gluon
PDFs \cite{Wang:2017eel,Wang:2017qyg,Wang:2019tgg}, for the
pion distribution amplitude (DA) \cite{Ji:2015qla} and generalized parton distributions (GPDs)
\cite{Ji:2015qla,Xiong:2015nua,Liu:2019urm}.

Within the pseudo-PDF approach,
the matching relations were derived for non-singlet
PDFs \cite{Ji:2017rah,Radyushkin:2017lvu,Radyushkin:2018cvn,Zhang:2018ggy,Izubuchi:2018srq}.
The strategy of the lattice extraction of non-singlet GPDs and the pion DA
using the pseudo-PDF methods was outlined
in a recent paper Ref. \cite{Radyushkin:2019owq},
where the matching conditions for these cases have been also derived.
In the present paper, our main goal is to describe the basic points
of the pseudo-PDF approach to extraction of unpolarized gluon PDFs, and also to find
one-loop matching conditions.

In the gluon case, the calculation is complicated by
strict requirements of gauge invariance. In this situation,
a very effective method is provided by the coordinate-representation approach of
Ref. \cite{Balitsky:1987bk}. It is based on the background-field method
and the heat-kernel expansion. It allows, starting with the original gauge-invariant
bilocal operator, to find its modification by one-loop corrections.
The results are obtained in an explicitly gauge-invariant form.

In this approach, there is no need
to specify the nature of matrix element characteristic of a particular
parton distribution. This means that one and the same Feynman diagram
calculation may be used both for finding
matching conditions for PDFs (given by forward matrix elements),
and for DA's and GPDs corresponding to non-forward ones (see Ref. \cite{Radyushkin:2019owq}
for an illustration of how this works for quark operators).

The paper is organized as follows.
In
\mbox{Section 2}, we analyze the kinematic structure of the matrix elements of the gluonic bilocal operators,
and identify those that contain information about the twist-2 gluon PDF.

Next, we discuss one-loop corrections. In \mbox{Section 3,} we analyze the gauge-link self-energy contribution
and specific properties of its ultraviolet and short-distance behavior.
Our results for the vertex corrections to the gluon link are given in \mbox{Section 4}
in the form that is valid both in forward and non-forward cases.
The ``box'' diagram is discussed in Section 5.
Since our results in this case are rather lengthy,
we present just some of them, and in the forward case only.
The gluon self-energy corrections are discussed in Section 6.

The subject of Section 7 is the structure of perturbative evolution
of the gluon operators and matching conditions.
Section 8 contains a summary of the paper.
